\documentclass[10pt,a4paper]{amsart}
\usepackage[margin=19mm]{geometry}
\usepackage{amsmath,amssymb,mathtools}
\usepackage[scr=rsfs,cal=euler]{mathalfa}
\usepackage{xfrac}
\usepackage[unicode,hidelinks,bookmarks=false]{hyperref}
\newcommand{\ie}{i.e.\ }
\newcommand{\cf}{cf.\ }
\newcommand{\resp}{respectively\ }
\newcommand{\Subsection}{\subsection}
\providecommand{\nobreakdash}{\nobreak\hbox{-}}
\setlength{\emergencystretch}{2em}
\multlinegap0pt
\usepackage{bm}
\usepackage{bbm}
\usepackage{dsfont}
\usepackage{xspace}
\usepackage{cancel}
\usepackage{mathtools}
\usepackage{amsthm}
\makeatletter
\@ifundefined{theorem}{\newtheorem{theorem}{Theorem}}{}
\makeatother
\newcommand{\Ric}{\text{Ric}}
\newcommand{\dive}{\text{div}}
\newcommand{\p}{\partial}
\newcommand{\sff}{\mathrm{I\!I}}
\title[Free Boundary Stable Minimal Hypersurfaces in Positively Cur]{Free Boundary Stable Minimal Hypersurfaces in Positively Curved 4-Manifolds}
\author{Yujie Wu}
\date{Source: arXiv:2308.08103v3 (CC BY 4.0)}
\begin{document}
\maketitle

\noindent\emph{Source and licence.} Free Boundary Stable Minimal Hypersurfaces in Positively Curved 4-Manifolds. Version arXiv:2308.08103v3 (CC BY 4.0), distributed under the Creative Commons Attribution 4.0 International licence, as recorded in the arXiv licence metadata for this exact version and shown on the abstract page https://arxiv.org/abs/2308.08103v3. Full rights notice: licenses/arxiv-2308.08103-CC-BY-4.0.txt.\par\smallskip

\noindent\textit{Editorial scope.} Counted unit: the Theorem whose source label is \texttt{VarForm} in the article (source0.tex lines 679--692, source label \texttt{VarForm}), delivered together with the article's own complete original proof of it (source0.tex lines 694--753, one \texttt{proof} environment of 60 lines). No mathematics is written, repaired or re-derived: statement and proof are the article's own text line for line. Required context retained uncounted: none. Every definition, display and label of those lines is retained unchanged. Omitted from this package and not redistributed: 1--678, 693--693, 754--1010. Nothing inside a retained excerpt is omitted or altered: every retained body is delivered exactly as the article's lines are written, so no omission receipt of any kind accompanies this package. Citations: the retained excerpts carry no citation command at all, so no bibliography entry of the article is transcribed and no reference is redistributed. Reference adaptation: none was needed and none was made; every reference inside the delivered unit resolves inside it, so no unresolved reference or citation is produced. Printed slips kept literal: no printed word, symbol, hypothesis or number of the counted statement or of its proof was altered. Layout and class: the article's own class is \texttt{amsart} with packages including amsmath, amssymb, hyperref, amsfonts, xypic, geometry, physics, fontenc, amsthm, graphicx, subcaption, euscript, multicol, xcolor; the wrapper is the lane's pinned \texttt{amsart} editorial wrapper at 10pt on A4, which restates the theorem-like environments the retained text needs and adds the article's own macro definitions that the retained text needs (\texttt{\textbackslash Ric}, \texttt{\textbackslash dive}, \texttt{\textbackslash p}, \texttt{\textbackslash sff}), with extra wrapper packages (bm, bbm, dsfont, xspace, cancel, mathtools). The delivered numbering is the unit's own. Assets: no retained span contains a figure environment, a graphics-inclusion command, a PDF-page inclusion, a table or a TikZ picture; no image file is copied into this package, the package declares no asset dependency and no image file is redistributed with it. Licence: version arXiv:2308.08103v3 of this article is distributed under the Creative Commons Attribution 4.0 International licence, as recorded in the arXiv licence metadata for this exact version and shown on the abstract page https://arxiv.org/abs/2308.08103v3. A full rights notice is delivered at licenses/arxiv-2308.08103-CC-BY-4.0.txt. Characters: every retained excerpt is pure ASCII, so the delivered engine, XeLaTeX with its default Latin Modern text font, reports no missing character.
	\begin{theorem} \label{VarForm}
		Assume $\Omega$ is a minimizer of $\mathcal{A}$ in the settings above, we have the following first  variation formula, writing $\Sigma =\p \Omega$,
		\begin{equation*}
			\nabla_{\nu_{\Sigma}} u-hu+uH_{\Sigma} =0 \text{ on } \Sigma, \quad \nu_{\partial \Sigma} (x) \perp T_x\partial N \quad \text{for }x \in \partial \Sigma \subset \partial_0 N ,
		\end{equation*}
		and the second variation formula,
		\begin{align*}
			&\frac{d^2}{dt^2}\Big\vert_{t=0}(\mathcal{A}(\varphi_t(\Omega)))\\
			=&\int_{\Sigma} |\nabla_{\Sigma} \phi|^2u  -\frac{u\phi^2}{2}(R_{N}-R_{\Sigma}+|\sff|^2+H^2_{\Sigma})+ \phi^2 (\Delta_N u-\Delta_{\Sigma}u-\nabla_{\nu_{\Sigma}} (hu))\\
			& -\int_{\partial \Sigma} u  \phi^2 A(\nu_{\Sigma},\nu_{\Sigma})\\
			\leq & \int_{\Sigma} |\nabla \phi|^2u  -\frac{u\phi^2}{2}(R_{N}-R_{\Sigma})+\phi^2(\Delta_N u-\Delta_{\Sigma}u-u\nabla_{\nu_{\Sigma}} h-\frac{h^2u}{2}-\frac{u^{-1}}{2}(\nabla_{\nu_{\Sigma}}u)^2)\\
			& -\int_{\partial \Sigma} u  \phi^2 A(\nu_{\Sigma},\nu_{\Sigma})
		\end{align*}
	\end{theorem}

	\begin{proof}
		The computation follows similarly from first and second variation formula of free boundary minimal hypersurfaces.
		We consider a family of diffeomorphism $\varphi_t$ of $N$ with vector field $X_t$, notice that if $x \in \partial N$ then $X_t \in T_x \partial N$. Let $\partial \Omega_t =:\Sigma_t$, the first variation is given as:
		\begin{align*}
			\frac{d}{dt}(\mathcal{A}(\varphi_t(\Omega)))&=\frac{d}{dt} \int_{\partial \Omega_t} u - \frac{d}{dt}\int_{\Omega_t} hu\\
			&=\int (\frac{d}{dt}u)dvol_{\partial \Omega_t}+ \int u \frac{d}{dt} dvol_{\partial_{\Omega_t}} - \int hu \frac{d}{dt} dvol_{\Omega_t}\\
			&= \int (\nabla_{X_t} u) dvol_{\partial \Omega_t} -\int (hu) \langle X_t, \nu_{\partial \Omega_t}\rangle dvol_{\partial \Omega_t} + \int (u \dive_{\partial \Omega_t} X_t) dvol_{\partial \Omega_t} \\
			&= \int (\nabla_{X_t} u) dvol_{\partial \Omega_t}  -\int (hu)\langle X_t, \nu_{\partial \Omega_t}\rangle  dvol_{\partial \Omega_t}+\int (u \dive_{\partial \Omega_t} X_t) dvol_{\partial \Omega_t}\\
			&=\int (\nabla_{X_t} u) dvol_{\Sigma_t}  -\int (hu) \langle X_t, \nu_{\Sigma_t}\rangle dvol_{\Sigma_t}+ \int (u \dive_{\Sigma_t} (X^{\perp}_t+X^{\top}_t)) dvol_{\Sigma_t}\\
			&=\int (\langle \nabla u, (X_t-X^{\top}_t)\rangle - hu\langle X_t, \nu_{\Sigma_t}\rangle-u\vec{H_{\Sigma_t}}\cdot X^{\perp}_t ) dvol_{\Sigma_t}+\int_{\partial \Sigma_t} u \langle X_t, \nu_{\partial \Sigma_t}\rangle \\
			&= \int (\nabla u \cdot X^{\perp}_t -hu\langle X_t, \nu_{\Sigma_t}\rangle-u\vec{H_{\Sigma_t}}\cdot X^{\perp}_t ) dvol_{\Sigma_t} + \int_{\partial \Sigma_t} u \langle X_t, \nu_{\partial \Sigma_t}\rangle
		\end{align*}

		Note we used the convention that mean curvature $H$ is defined as the trace of the second fundamental form and hence $H_{\Sigma}=-\langle\nabla_{e_i} e_i,\nu_{\Sigma}\rangle=-\vec{H_{\Sigma}}\cdot \nu_{\Sigma}$.
		So at $t=0$ we have $\nabla_{\nu_{\Sigma}} u-hu+uH_{\Sigma} =0$ on $\Sigma,$ and that $\nu_{\partial \Sigma} (x) \perp T_x\partial N$ for $x \in \partial \Sigma \subset \partial N.$

		Now we continue with the second variation :
		\begin{align*}
			&\frac{d}{dt}\Big\vert_{t=0}(\mathcal{A}'(\varphi_t(\Omega))-\int_{\partial \Sigma_t} u \langle X_t, \nu_{\partial \Sigma_t}\rangle)\\
			=& \int \frac{d}{dt}\Big\vert_{t=0}(\nabla u \cdot X^{\perp}_t -hu\langle X_t, \nu_{\Sigma_t}\rangle-u\vec{H_{\Sigma_t}}\cdot X^{\perp}_t ) dvol_{\Sigma}\\
			=&\int \phi_t \frac{d}{dt} \Big\vert_{t=0}(\nabla u \cdot \nu_{\Sigma_t} -hu+uH_{\Sigma_t})dvol_{\Sigma}\\
			=& \int \phi_t (\partial_t \langle\nabla u,  \nu_{\Sigma_t}\rangle-\nabla_{X_t}(hu)+ (\nabla_{X_t} u) H_{\Sigma_t}+ u \partial_t H_{\Sigma_t}) dvol_{\Sigma},  \quad \text{at $t=0$}.\\
		\end{align*}
			Since at $t=0, \partial \Sigma \cap \partial N$ orthogonally, using the exponential map near $\partial \Sigma$, for any smooth function $\phi_t$, the diffeomorphism near $\Sigma$ given by $\Sigma \times (-\epsilon,\epsilon) \ni (x,t) \rightarrow \exp_x(t\nu_x)$ is admissible and produce a normal variation near $\Sigma$.
			We will also use $\Delta_N u-\Delta_{\Sigma}u=\nabla^2 u(\nu_{\Sigma},\nu_{\Sigma})+H_{\Sigma}\nabla_{\nu_{\Sigma}}u$:
		\begin{align*}
			&\frac{d}{dt}\Big\vert_{t=0}(\mathcal{A}'(\varphi_t(\Omega))-\int_{\partial \Sigma_t} u \langle X_t, \nu_{\partial \Sigma_t}\rangle)\\
			=&\frac{d}{dt}\Big\vert_{t=0}(\mathcal{A}'(\varphi_t(\Omega)))\\
			=& \int \phi_t^2(\nabla^2u(\nu_{\Sigma_t},\nu_{\Sigma_t})-\nabla_{ \nu_{\Sigma_t}}(hu)+H_{\Sigma_t}\nabla_{ \nu_{\Sigma_t}}u) +\phi_t \langle\nabla u,\p_t \nu_{\Sigma_t}\rangle dvol_{\Sigma}\\
			&+ \int u \phi_t(-\Delta_{\Sigma_t} \phi_t-\phi_t(|\sff_{\Sigma_t}|^2+\Ric_N(\nu_{\Sigma_t},\nu_{\Sigma_t})))dvol_{\Sigma}  ,  \quad \text{at $t=0$} \\
			=& \int_{\Sigma} |\nabla_{\Sigma} \phi|^2u  -u\phi^2(|\sff_{\Sigma}|^2+\Ric_N(\nu_{\Sigma},\nu_{\Sigma}))+ \phi^2 \nabla^2 u(\nu_{\Sigma},\nu_{\Sigma})+\phi_t \langle\nabla u,\p_t \nu_{\Sigma_t}\rangle\\
			&-\int_{\Sigma} \phi^2(\nabla_{ \nu_{\Sigma}}(hu)-H_{\Sigma}\nabla_{ \nu_{\Sigma}}u)+\int_{\Sigma} \langle \nabla_{\Sigma} u, \nabla_{\Sigma} \phi \rangle \phi   -\int_{\partial \Sigma} u\phi \nabla_{\nu_{\partial \Sigma}} \phi\\
			=&\int_{\Sigma} |\nabla_{\Sigma} \phi|^2u  -u\phi^2(|\sff_{\Sigma}|^2+\Ric_N(\nu_{\Sigma},\nu_{\Sigma}))+ \phi^2 (\Delta_N u-\Delta_{\Sigma}u)-\phi^2\nabla_{\nu_{\Sigma}} (hu)\\
			&+ \int_{\Sigma} \phi_t \langle\nabla u,\p_t \nu_{\Sigma_t}\rangle +\langle \nabla_{\Sigma} u, \nabla_{\Sigma} \phi \rangle \phi   -\int_{\partial \Sigma} u\phi \nabla_{\nu_{\partial \Sigma}} \phi
		\end{align*}
		Now we use that for a family of evolution of hypersurfaces using a normal vector field we have $-\nabla_{\Sigma} \phi=\partial_t \nu_{\Sigma}$ and from the Gauss equation we have $R_N=R_{\Sigma}+2\Ric(\nu,\nu)+|\sff|^2-H^2_{\Sigma}$ to get,
		\begin{align*}
			&\frac{d^2}{dt^2}\Big\vert_{t=0}(\mathcal{A}(\varphi_t(\Omega)))\\
			=& \int_{\Sigma} |\nabla_{\Sigma} \phi|^2u  -\frac{u\phi^2}{2}(R_{N}-R_{\Sigma}+|\sff|^2+H^2_{\Sigma})+ \phi^2 (\Delta_N u-\Delta_{\Sigma}u-\nabla_{\nu_{\Sigma}} (hu))\\
			&-\int_{\partial \Sigma} u\phi \langle \nabla_{\Sigma}\phi, \nu_{\partial \Sigma}\rangle\\
			=&\int_{\Sigma} |\nabla_{\Sigma} \phi|^2u  -\frac{u\phi^2}{2}(R_{N}-R_{\Sigma}+|\sff|^2+H^2_{\Sigma})+ \phi^2 (\Delta_N u-\Delta_{\Sigma}u-\nabla_{\nu_{\Sigma}} (hu))\\
			&+ \int_{\partial \Sigma} u\phi \langle \partial_t \nu_{\Sigma}, \nu_{\p \Sigma} \rangle ,\quad \text{at } t=0, \\
			=&\int_{\Sigma} |\nabla_{\Sigma} \phi|^2u  -\frac{u\phi^2}{2}(R_{N}-R_{\Sigma}+|\sff|^2+H^2_{\Sigma})+ \phi^2 (\Delta_N u-\Delta_{\Sigma}u-\nabla_{\nu_{\Sigma}} (hu))\\
			& + \int_{\partial \Sigma} u  \langle \nabla_{X_t} X_t, \nu_{\p \Sigma} \rangle ,\quad \text{at } t=0, \\
			=&\int_{\Sigma} |\nabla_{\Sigma} \phi|^2u  -\frac{u\phi^2}{2}(R_{N}-R_{\Sigma}+|\sff|^2+H^2_{\Sigma})+ \phi^2 (\Delta_N u-\Delta_{\Sigma}u-\nabla_{\nu_{\Sigma}} (hu))\\
			& - \int_{\partial \Sigma} u  \phi^2 A(\nu_{\Sigma},\nu_{\Sigma})
		\end{align*}
		
		We now write $|\sff|^2=|\mathring{\sff}|^2+\frac{H^2_{\Sigma}}{2}\geq  \frac{H^2_{\Sigma}}{2}$ 
		and notice that according to the first variation and $u>0$ we have $\frac{u^{-1}}{2}(uH_{\Sigma})^2=\frac{u^{-1}}{2}(\nabla_{\nu_{\Sigma}}u)^2+\frac{h^2u}{2}-h\nabla_{\nu_{\Sigma}}u$, so in total:

		\begin{align*}
			0\leq &\frac{d^2}{dt^2}\Big\vert_{t=0}(\mathcal{A}(\varphi_t(\Omega)))\\
			\leq& \int_{\Sigma} |\nabla \phi|^2u  -\frac{u\phi^2}{2}(R_{N}-R_{\Sigma})+\phi^2(-\frac{3H^2_{\Sigma}}{4}u+ \Delta_N u-\Delta_{\Sigma}u-\nabla_{\nu_{\Sigma}} (hu))\\
			& -\int_{\partial \Sigma} u  \phi^2 A(\nu_{\Sigma},\nu_{\Sigma})\\
			\leq & \int_{\Sigma} |\nabla \phi|^2u  -\frac{u\phi^2}{2}(R_{N}-R_{\Sigma})+\phi^2(\Delta_N u-\Delta_{\Sigma}u-u\nabla_{\nu_{\Sigma}} h-\frac{h^2u}{2}-\frac{u^{-1}}{2}(\nabla_{\nu_{\Sigma}}u)^2)\\
			& -\int_{\partial \Sigma} u  \phi^2 A(\nu_{\Sigma},\nu_{\Sigma})\\
		\end{align*}
		
	\end{proof}

\end{document}
